\chapter{Coarse geometry of trees}
\label{ch:geometry}

The two obstructions below are deterministic statements about arbitrary infinite trees. They
are proved over Mathlib's \texttt{SimpleGraph} in its graph metric, and no hyperbolicity
enters.

\begin{definition}[Quasi-isometry]
\label{def:qi}
\lean{BranchingProcess.IsQIWith, BranchingProcess.QuasiIsometric}
\leanok
A map $f\colon V\to V'$ is a \emph{$D$-quasi-isometry} from $G$ to $G'$ when
\[
  d'(f x,f y)\leq D\,d(x,y)+D,
  \qquad
  d(x,y)\leq D\,d'(f x,f y)+D^2,
\]
for all $x,y$, and every vertex of $G'$ lies within $D$ of the image. Two graphs are
\emph{quasi-isometric} when such a map exists for some $D$.
\end{definition}

\begin{definition}[Betweenness, rays, lines and hairs]
\label{def:between}
\lean{BranchingProcess.Between, BranchingProcess.IsRay, BranchingProcess.IsLine,
BranchingProcess.IsHair, BranchingProcess.rayGraph}
\leanok
A vertex $m$ lies \emph{between} $u$ and $v$ when every walk from $u$ to $v$ meets $m$. A
\emph{ray} is a map $\bbN\to V$ isometric onto its image and a \emph{line} is such a map
$\bbZ\to V$. A \emph{hair} at $c$ of depth $h$ is a component $P$ of the complement of
$\set{c}$ with $d(v,c)\leq h$ for every $v\in P$ and equality for some $v$. Finally
$\mathrm{Ray}$ is $\bbN$ with the path-graph structure.
\end{definition}

\begin{lemma}[Quasi-geodesic stability in a tree]
\label{lem:stability}
\uses{def:qi, def:between}
\lean{BranchingProcess.between_of_mem_support, BranchingProcess.between_iff_dist_add,
BranchingProcess.exists_between_dist_le}
\leanok
Let $f$ be a $D$-quasi-isometry between trees. Every vertex separating $f(u)$ from $f(v)$ lies
within $D$ of the image of a walk from $u$ to $v$.
\end{lemma}

\begin{proof}
\leanok
Let $p$ be a walk from $u$ to $v$, and write $p_i$ for its $i$-th vertex. Suppose for a
contradiction that every $f(p_i)$ is more than $D$ from $m$. We show inductively that $m$
does not lie between $f(u)$ and $f(p_i)$. This is clear for $i=0$. If it fails for $i+1$,
the tree property applied at $f(p_i)$ shows that $m$ lies either between $f(u)$ and
$f(p_i)$ or between $f(p_i)$ and $f(p_{i+1})$. The first alternative contradicts the
induction hypothesis. In the second, metric additivity and the assumed lower bounds give
\[
  d'(f(p_i),f(p_{i+1}))>2D.
\]
However, $p_i$ and $p_{i+1}$ are adjacent, so the upper quasi-isometry bound gives
$d'(f(p_i),f(p_{i+1}))\leq2D$. At $i=n$ the induction contradicts the assumption that $m$
lies between $f(u)$ and $f(v)$. Hence some $f(p_i)$ lies within $D$ of $m$.
\end{proof}

\begin{theorem}[Hairs against lines]
\label{thm:hair-separation}
\uses{lem:stability}
\lean{BranchingProcess.not_quasiIsometric_of_hair_of_line}
\leanok
A tree carrying hairs of unbounded depth is not quasi-isometric to a tree all of whose
vertices lie within a bounded distance of a line.
\end{theorem}

\begin{proof}
\leanok
Suppose that $f$ is a $D$-quasi-isometry, and choose a hair $P$ at $c$ of depth
\[
  h>D(D+R)+D^2.
\]
Let $x\in P$ satisfy $d(x,c)=h$. Choose a line $l$ and $k\in\bbZ$ such that
$d'(f(x),l(k))\leq R$. One of the two directions along $l$ contains, arbitrarily far from
$l(k)$, a vertex $z$ for which $l(k)$ lies between $f(c)$ and $z$. Coarse surjectivity
provides $p$ with $d'(f(p),z)\leq D$. Taking $z$ sufficiently far away ensures that
$p\notin P$, since every point of $P$ is within $h$ of $c$ and the upper quasi-isometry
bound would otherwise keep $f(p)$ too close to $f(c)$.

The vertex $l(k)$ lies between $f(c)$ and $f(p)$. Apply \cref{lem:stability} to a geodesic
from $c$ to $p$ to obtain a vertex $y$ on that geodesic with
$d'(f(y),l(k))\leq D$. Since $x\in P$ and $p\notin P$, the vertex $c$ separates $p$ from
$x$. As $y$ lies between $c$ and $p$, it follows that $c$ also lies between $y$ and $x$,
and hence $d(y,x)\geq d(c,x)=h$. On the other hand,
$d'(f(y),f(x))\leq D+R$, so the lower quasi-isometry bound gives
\[
  d(y,x)\leq D(D+R)+D^2<h,
\]
a contradiction.
\end{proof}

\begin{theorem}[Three rays]
\label{thm:three-rays}
\uses{lem:stability}
\lean{BranchingProcess.not_quasiIsometric_rayGraph, BranchingProcess.rayGraph_isTree,
BranchingProcess.dist_ray_branch}
\leanok
A tree carrying three rays that pairwise meet only in their common initial vertex is not
quasi-isometric to $\mathrm{Ray}$.
\end{theorem}

\begin{proof}
\leanok
Suppose that $f$ is a $D$-quasi-isometry and set $\rho=2D^2+1$. For distinct rays $r_i$
and $r_j$, uniqueness of paths in the source tree gives
\[
  d(r_i(m),r_j(n))=m+n.
\]
Among the three vertices $f(r_i(\rho))$ of $\mathrm{Ray}$, two lie on the same side of
$f(v)$. After exchanging them if necessary, $f(r_i(\rho))$ lies between $f(v)$ and
$f(r_j(\rho))$.

Apply \cref{lem:stability} to the geodesic from $v$ to $r_j(\rho)$. It gives a vertex $y$
on this geodesic such that
$d(f(y),f(r_i(\rho)))\leq D$. The geodesic is the initial segment of $r_j$, so
$y=r_j(t)$ for some $t\leq\rho$. The two rays meet only at $v$, and consequently
\[
  d(y,r_i(\rho))=t+\rho\geq\rho.
\]
The lower quasi-isometry bound gives the opposite estimate
\[
  d(y,r_i(\rho))\leq D\,d(f(y),f(r_i(\rho)))+D^2
  \leq2D^2<\rho,
\]
which is the required contradiction.
\end{proof}
